%=====================================================================
\chapter{Python Reference for Artificial Intelligence}\label{app:python}
%=====================================================================
\normalfigures

Everything the labs use, on a few pages. Each lab is also a runnable file in the
\texttt{code/} folder, named after its chapter.

\section{Setting Up}

Any Python 3.10 or later will do. Chapters~1--2, 4--5 and~7--13 and~17--22 use the
standard library only; Chapters~3, 6, 15 and~16 also use NumPy. Check with:
\begin{lstlisting}[language=Python]
import sys, numpy
print(sys.version)
print(numpy.__version__)
\end{lstlisting}
The labs were tested with Python 3.12 and NumPy 2.3. Run one from a terminal with
\texttt{python ch05\_over\_sprint\_states\_two.py}, or paste it into a notebook
cell. No lab needs a GPU, a network connection or an account with anybody.

\section{The Four Structures That Matter}

\rowcolors{2}{white}{lightbg}
\begin{longtable}{@{}L{6.0cm}L{9.0cm}@{}}
\toprule
\rowcolor{primary!15}
\textbf{Code} & \textbf{What it does} \\
\midrule
\endhead
\texttt{frozenset(\{"a", "b"\})} & An immutable, hashable set. The state
representation of Chapters 3--5 and 9 \\
\texttt{s | \{x\}} & A \emph{new} set with \texttt{x} added. Never
\texttt{s.add(x)} on a state \\
\texttt{deps <= done} & Subset test: every dependency is finished \\
\texttt{a \& b} & Intersection --- non-empty means ``at least one of these is
done'' \\
\texttt{Node = namedtuple("Node", "state parent action cost")} & A tuple with
named fields: hashable, comparable, self-documenting when printed \\
\texttt{("S1", None, "S3")} & A tuple of slots: a partial assignment, with
\texttt{None} where a variable is unassigned (Chapters 7, 11) \\
\texttt{d = \{\}; d.setdefault(k, []).append(v)} & Build an adjacency dictionary
without checking whether the key exists \\
\bottomrule
\end{longtable}

\section{Frontiers}

\begin{lstlisting}[language=Python]
from collections import deque
import heapq, itertools

q = deque()          # breadth-first
q.append(node)
node = q.popleft()   # O(1). NEVER list.pop(0), which is O(n)

st = []              # depth-first
st.append(node)
node = st.pop()

h, tie = [], itertools.count()          # best-first / A*
heapq.heappush(h, (priority, next(tie), node))
node = heapq.heappop(h)[2]
\end{lstlisting}
The counter in the middle of the heap entry breaks ties by insertion order, so two
nodes are never compared. Omitting it produces an intermittent
\texttt{TypeError} --- see \chref{python}.

\section{Generators, Memoisation and Recursion}

\begin{lstlisting}[language=Python]
def actions(state):              # a successor function
    for t, (_, deps) in TASKS.items():
        if t not in state and deps <= state:
            yield t              # produced one at a time, on demand

from functools import lru_cache

@lru_cache(maxsize=None)         # arguments must be hashable
def minimax(state, maximising):
    ...

import sys
sys.setrecursionlimit(10000)     # deep proof trees, Chapters 11-12
\end{lstlisting}

\section{Exact Arithmetic for Probability}

Chapters~14 and~15 print probabilities that you should be able to check by hand.
\texttt{Fraction} keeps them exact:
\begin{lstlisting}[language=Python]
from fractions import Fraction as F

p = F(1, 3) + F(1, 6)        # Fraction(1, 2) -- exactly
print(p, float(p))           # 1/2 0.5
\end{lstlisting}

\section{NumPy, Where It Earns Its Place}

Only four chapters use it, and only for grids of numbers.
\begin{lstlisting}[language=Python]
import numpy as np

x = np.linspace(0, 10, 101)          # 101 points, Chapter 16 memberships
mu = np.clip((x - 2) / 3, 0, 1)      # a shoulder membership function
area = np.trapezoid(mu, x)           # centroid numerator/denominator
rng = np.random.default_rng(0)       # seeded sampler, Chapter 15
s = rng.random(10000)
\end{lstlisting}

\section{Reproducibility}

Every lab that samples seeds its generator explicitly, with
\texttt{random.Random(0)} or with NumPy's \texttt{default\_rng(0)}, so the numbers
printed in this book are the numbers you will see. If your output differs, check
the seed before checking the algorithm.

\section{Style Rules the Labs Follow}

\begin{tight}
  \item \textbf{Pure ASCII.} Write \texttt{AND}, \texttt{NOT}, \texttt{->},
        \texttt{forall} --- never the Unicode symbols. \texttt{code/sync\_labs.py}
        rejects any byte above 127, because the handout's listings are typeset by
        pdf\LaTeX, which cannot set them.
  \item \textbf{78 columns.} Longer lines are wrapped by the same check.
  \item \textbf{The handout is the source.} \texttt{sync\_labs.py} extracts each
        lab from its chapter into \texttt{code/}. Edit the chapter, not the
        extracted file.
  \item \textbf{Assertions, not commentary.} Each lab ends by asserting the claims
        the chapter makes about it, so a broken claim fails loudly.
\end{tight}

\begin{lstlisting}
cd code
python sync_labs.py          # extract every chapter's lab
python sync_labs.py --run    # ... and run them all
python sync_labs.py --run 9  # ... run chapter 9 only
\end{lstlisting}
